%% kapa: exam=1 round=23 year=2012 period=1 form=A subject=economics subject_ko=경제원론 seq=37 number=77
%% source: Q-Net 감정평가사 기출문제 (cst00309) artlSeq=5052578 file=raw/1008/5052578/2012년 제23회 감정평가사 제1차시험 기출문제(1교시 A형).pdf
%% conversion: kordoc 4.24.51 -> tools/gen1.py (auto); review=proofread-sonnet; flags=none
\begin{question}{77}
중앙은행은 아래와 같은 테일러 준칙(Taylor rule)에 따라 명목이자율을 조정한다. 이에 관한 설명으로 옳지 않은 것은? (단, $i$는 명목이자율, $\pi$는 인플레이션율, $\pi^{*}$는 목표 인플레이션율, $Y^{*}$는 잠재 GDP, $Y$는 실제 GDP, $(Y^{*}-Y)/Y^{*}$는 총생산 갭이다.)
\begin{framedtext}
\centering $i=0.05+\pi+0.5(\pi-\pi^{*})-0.5(Y^{*}-Y)/Y^{*}$
\end{framedtext}
\begin{choices}
  \item 목표 인플레이션율이 낮아지면 중앙은행은 명목이자율을 인상한다.
  \item 실제 GDP가 잠재 GDP보다 더 큰 경우에 중앙은행은 명목이자율을 인상한다.
  \item 총생산 갭은 0이고 인플레이션율이 3\%에서 4\%로 상승하는 경우에, 중앙은행은 명목이자율을 0.5\%포인트(\%p) 인상한다.
  \item 인플레이션율이 목표치와 같고 실제 GDP가 잠재 GDP와 같다면 실질이자율은 5\%가 된다.
  \item 인플레이션율은 목표치와 같고 총생산 갭이 0\%에서 1\%로 상승하는 경우에, 중앙은행은 명목이자율을 0.5\%포인트(\%p) 인하한다.
\end{choices}
\end{question}
\answer{3}
